%% RNAAS skeleton: Gaia DR4 constraints on outer companions to transiting giant-planet hosts.
%% Written 2026-10-12, before DR4. Upload to Overleaf with detections.tex (scripts/make_table.py)
%% and references.bib.
%%
%% Rules (notes/aas_ai_policy.md, CLAUDE.md sections 5-6):
%%  - Word budget: the submission site states 1,350 words + a 150-word abstract (the guidelines
%%    page says 1,500 incl. references); plan for 1,350 body words.
%%  - ONE table (detections.tex); no figure. Full per-host table as MRT (hosts.mrt).
%%  - Single author: write "I" in the body. No line numbers.
%%  - Every number traces to a file in results/ and the script that wrote it; mark the source
%%    in a comment next to each number.
%%  - \XX{...} marks a value filled in after release. None may remain at submission.
%%  - References: export BibTeX from ADS, or verify each entry against ADS/arXiv by hand.
%%    Never unverified entries.
%%  - Matches the submitted duotransits note (AAS82207; C:\git\duo-transits,
%%    singles/rnaas/rnaas.tex at d932e46): Overleaf AAS template "AASTeX Template for
%%    submissions to AAS Journals (ApJ-AJ-ApJS-ApJL-PSJ-RNAAS)" (aastex701.cls,
%%    aasjournalv7.bst, orcid-ID.png), plain \documentclass{aastex701} (the RNAAS option is
%%    deprecated in v7). Submission-site steps: notes/aas_ai_policy.md.
\documentclass{aastex701}

\usepackage{xcolor}  % drafting only (for \XX); remove with the last \XX before submission
\newcommand{\XX}[1]{\textcolor{red}{[#1]}}
\newcommand{\gaia}{\textit{Gaia}}

\begin{document}

\title{\gaia\ DR4 Constraints on Outer Giant Companions to Stars Hosting Transiting Giant Planets}

\author[orcid=0009-0000-7524-4470]{Jon McBee}
\affiliation{Independent researcher, USA}
\email{jon.mcbee@gmail.com}

\begin{abstract}
% <= 150 words. Structure: what DR4 enables; sample; method in one sentence; detections;
% occurrence / limits; mutual inclinations if any; data release.
\gaia\ Data Release 4 publishes per-transit epoch astrometry for the first time. I searched it
for astrometric companions to \XX{N_hosts} stars hosting transiting giant planets, fitting a
single-star model, two acceleration models and a Keplerian orbit (periods 0.1--4\,yr) with
maximum-likelihood excess noise, at a false-alarm probability of $10^{-4}$ per star calibrated on
simulated DR4 data. I find \XX{N_orbit} orbits and \XX{N_acc} accelerations; \XX{summary of
vetting}. For companions of 1--13\,$M_J$ at 1--4\,yr, the occurrence rate is \XX{f} (\XX{95\%
upper limit}). \XX{Mutual-inclination sentence if detections allow.} The per-host results and
completeness limits are released as a machine-readable table.
\end{abstract}

% UAT concepts with their IDs from https://astrothesaurus.org (verify each, e.g. astrometric
% exoplanet detection, hot Jupiters, exoplanet systems). Entered again on the submission site.
\keywords{\XX{UAT concepts}}

\section{Introduction}
% ~150 words. Hot/warm Jupiter migration; outer perturbers; RV companion rates (Knutson et al.
% 2014; Bryan et al. 2016); HGCA; what astrometry adds (true mass, inclination; transit fixes
% i_b ~ 90 deg, so minimum mutual inclinations). Xiao et al. 2026. Lammers & Winn forecasts.
\XX{Introduction.}

\section{Data and method} \label{sec:method}
% ~400 words. Facts fixed before release; numbers below trace to results/ (see comments).
\paragraph{Sample}
Hosts of transiting planets with $M_p$ or $R_p$ in the giant range from the NASA Exoplanet
Archive, cross-matched to \gaia\ (\XX{N_hosts}; data/targets.csv, scripts/build_targets.py).

\paragraph{Model}
For each host I fit the \gaia\ DR4 along-scan positions (\texttt{centroid\_pos\_al}) of each CCD
with the six-parameter single-star model of the DR4 core solution (my fit reproduces ESA's
\texttt{gaiasupdate} to $10^{-9}$; notes/phase1\_validation.md), plus per-CCD and per-transit
excess noise fitted by maximum likelihood. Alternative models add a 7- or 9-parameter
acceleration or a Keplerian photocentre orbit (linear Thiele--Innes constants; grid search in
period, eccentricity and periastron time, then non-linear refinement). The statistic is twice
the log-likelihood gain over the single-star model.

\paragraph{Calibration}
Thresholds at a false-alarm probability of $10^{-4}$ per star come from $10^4$ simulated
single stars with the DR4 scanning law and noise model: \XX{42.6} (orbit), \XX{21.7} and
\XX{27.1} (accelerations) % results/thresholds.csv; scripts/calibrate_null.py
\XX{; confirmed / raised by \gaia\ DR4 control stars, results/thresholds\_real.csv}.
Injection--recovery of 6000 planets gives 50\% completeness at an expected full-orbit
$\Delta\chi^2 \approx 80$ % results/completeness_models.csv; scripts/injection_recovery.py
and per-host mass limits (\XX{results/target\_sensitivity.csv, recalibrated after release?}).

\paragraph{Vetting}
Each detection is checked for stability under removal of a fifth of the data, for unresolved
near-equal binaries (literature RVs, \gaia\ DR4 epoch RVs and double-lined products), for
known stellar companions (RUWE, Hipparcos--\gaia\ accelerations, DR3 non-single-star
solutions), against \gaia\ DPAC solutions, and by an independent re-derivation
(results/dr4/manual\_vetting.csv).

\section{Results} \label{sec:results}
% ~350 words + Table 1. Detections; accelerations; comparison with DPAC exoplanet catalogue
% (cite, don't claim); occurrence (results/dr4/population_summary.json); limits.
\XX{Results.}

\input{detections.tex}

\section{Discussion}
% ~150 words. Comparison with RV companion rates; mutual inclinations (minimum Delta I only for
% detections); caveats (impostors, P > 4 yr -> accelerations, DR4 noise).
\XX{Discussion.}

\begin{acknowledgments}
This work has made use of data from the European Space Agency (ESA) mission \gaia\
(\url{https://www.cosmos.esa.int/gaia}), processed by the \gaia\ Data Processing and Analysis
Consortium (DPAC, \url{https://www.cosmos.esa.int/web/gaia/dpac/consortium}). Funding for the
DPAC has been provided by national institutions, in particular the institutions participating
in the \gaia\ Multilateral Agreement. This research has made use of the NASA Exoplanet Archive,
which is operated by the California Institute of Technology, under contract with the National
Aeronautics and Space Administration under the Exoplanet Exploration Program.
% XX: check both texts against the current wording requested by ESA and the NASA Exoplanet
% Archive; shorten if the word budget is tight (the duotransits note used one sentence).
Data: \dataset[NASA Exoplanet Archive PSCompPars]{\doi{10.26133/NEA13}}, \gaia\ DR4 \XX{DOI}.
% NEA DOI verified on DataCite by the duotransits session; the submission site asks for
% dataset DOIs.
\end{acknowledgments}

\facilities{Gaia, Exoplanet Archive}
\software{astropy, astroquery, numpy, scipy, emcee, matplotlib, gaiasupdate, gaiascanlaw,
pandas \XX{versions; cite each via ADS / Zenodo / ASCL; code DOI for github.com/coracle-us/gaia-dr4}}

\section*{AI disclosure statement}
% Required by the AAS guidelines (notes/aas_ai_policy.md); placed after facilities/software and
% before the bibliography, worded as in the submitted duotransits note. Keep it TRUE for the
% final Note: edit if the division of work changes.
The author used an AI assistant, Anthropic's Claude (a large language model), under his
supervision. It wrote and ran the analysis code (\url{https://github.com/coracle-us/gaia-dr4})
and drafted and edited the text of this note. The author reviewed the validation tests
(reproduction of ESA's single-star fits, recovery of \gaia-4b, null and injection--recovery
calibration), reviewed and edited the text, checked the numbers against the saved analysis
outputs, and takes full responsibility for the content. Every reference was checked against
ADS records. The note contains no AI-generated figures.

\bibliography{references}{}
\bibliographystyle{aasjournalv7}

\end{document}
